\chapter{Software e dipendenze}

\section{Strumenti}
\begin{itemize}
  \item Arduino IDE o arduino-cli per compilazione e upload.
  \item CMake >= 3.16 per i test.
  \item Compilatore C++17 per i test su PC.
\end{itemize}
L'uso di arduino-cli consente automazione e integrazione con pipeline di test.

\section{Librerie}
\begin{itemize}
  \item Servo (Arduino)
  \item LiquidCrystal (Arduino)
  \item DHT (per DHT22)
  \item GoogleTest (solo per test)
\end{itemize}
Le librerie Arduino sono usate tramite header standard, mentre GoogleTest e scaricato dinamicamente nella fase di configurazione CMake.

\section{Repository}
Il progetto e organizzato in:
\begin{itemize}
  \item \texttt{Pinacoteca.ino} per il firmware;
  \item \texttt{lib/} per i moduli;
  \item \texttt{test/} per i test unitari;
  \item \texttt{diagram.json} per la simulazione Wokwi.
\end{itemize}
